% Direct-authority diplomatic French for printed p.157.
\SourcePage{157}

\LayerTarget{tohoku:def:sections-and-supports}
\noindent tel faisceau, $A$ une partie de $X$, on désigne par $\Gamma(A,F)$ le
groupe des sections de $F$ (considéré comme espace étalé) sur $A$, et on pose
$\Gamma(F)=\Gamma(X,F)$, (de sorte qu'on a
$\Gamma(A,F)=\Gamma(F\vert A)$, où $F\vert A$ désigne la restriction de $F$ à
$A$). Plus généralement, soit $\Phi$ un ensemble non vide de parties fermées
de $X$, filtrant croissant, et tel que $A\in\Phi$, $B\subset A$ implique
$B\in\Phi$ : pour abréger, nous dirons que $\Phi$ est un \textit{antifiltre
de parties fermées de $X$}. On désigne par $\Gamma_\Phi(F)$ le sous-groupe de
$\Gamma(F)$ formé des sections $f$ dont le support (= complémentaire du plus
grand ouvert dans $X$ sur lequel la restriction de $f$ soit nulle) est
$\in\Phi$.

\LayerTarget{tohoku:def:cohomology-with-supports}
On voit aussitôt que $\Gamma_\Phi(F)$ est un foncteur exact à gauche défini
sur $C^X$ à valeurs dans la catégorie des groupes abéliens. Il s'impose donc
de considérer ses foncteurs dérivés droits, identiques aux satellites droits
(cf. 2.3), qui existent en vertu du corollaire à la prop. 3.1.1. On les note
$H^p_\Phi(X,F)$, (où $p$ est un entier quelconque). D'après la théorie des
foncteurs dérivés :

\LayerTarget{tohoku:prop:3.2.1}
\textsc{Proposition 3.2.1.} --- \textit{Le système des foncteurs
$H^p_\Phi(X,F)$ ($-\infty<p<-\infty$) est caractérisé par les conditions
suivantes : ils forment un foncteur cohomologique sur $C^X$ à valeurs dans la
catégorie des groupes abéliens, $H^p_\Phi(X,F)$ est nul pour $p<0$, coïncide
avec $\Gamma_\Phi(F)$ pour $p=0$, et est effaçable (c'est-à-dire s'annule pour
$F$ injectif) si $p>0$.}

\LayerTarget{tohoku:construction:compute-cohomology-with-supports}
Pour calculer les $H^p_\Phi(X,F)$, on prend une résolution droite de $F$ par
des faisceaux $C^i$ injectifs, ou plus généralement tels qu'on sache à
l'avance que les $H^p_\Phi(X,C^i)$ sont nuls pour $p>0$ (on dit alors que les
$C^i$ sont $\Gamma_\Phi$-acycliques), on considère le complexe $C(F)$ formé
par les $C^i$, et on aura alors
\[
H^p_\Phi(X,F)=H^p\Gamma_\Phi(C(F)).
\]
Pour appliquer cette méthode, il est donc important de connaître des critères
permettant d'affirmer qu'un faisceau est $\Gamma_\Phi$-acyclique, nous en
indiquerons plus bas (cf. 3.3).

\LayerTarget{tohoku:construction:inverse-image-comparison:p1}
Soit $f$ une application continue d'un espace $X$ dans un espace $Y$, pour
tout faisceau $F$ sur $Y$, on définit de façon naturelle le \textit{faisceau
image réciproque} de $F$ par $f$, (noté $f^{-1}(F)$ par abus de notations) :
si on considère $F$ comme espace étalé sur $Y$, il suffit de reprendre la
définition de l'image inverse d'un fibré. On obtient ainsi un \textit{foncteur
covariant additif exact} $F\longmapsto f^{-1}(F)$ de $C^Y$ dans $C^X$, que
nous allons désigner par $g$. Pour tout $F\in C^Y$, on a un homomorphisme
évident $\Gamma(F)\longrightarrow\Gamma(g(F))$, d'ailleurs fonctoriel, i. e.
on a un homomorphisme de foncteurs
$\Gamma^Y\longrightarrow\Gamma^Xg$ (où nous avons mis en exposant l'espace
relativement auquel on considère le foncteur $\Gamma$). Plus généralement,
considérons un antifiltre $\Phi$ de parties fermées de $X$ et un antifilre
$\Psi$ de partie fermées de $Y$ tel que pour tout $B\in\Psi$, on ait
$f^{-1}(B)\in\Phi$. Alors pour tout $F\in C^Y$, l'homomorphisme
$\Gamma(F)\longrightarrow\Gamma(g(F))$ applique $\Gamma_\Psi(F)$ dans
$\Gamma_\Phi(g(F))$, d'où un homomorphisme fonctoriel
$\Gamma^Y_\Psi\longrightarrow\Gamma^X_\Phi g$. D'ailleurs, $g$ étant un
foncteur exact, le foncteur $(R\Gamma^X_\Phi)g$ peut être considéré comme un
foncteur cohomologique sur $C^Y$, se réduisant à $\Gamma^X_\Phi g$ en
dimensions 0 ; comme le foncteur cohomologique $R\Gamma^Y_\Psi$ est
\og universel\fg{} (prop. 2.2.1) l'homomorphisme fonctoriel
$\Gamma^Y_\Psi\longrightarrow\Gamma^X_\Phi g$

\TohokuContinuation{158}

